\documentclass[11pt,a4paper]{article}
\usepackage{pinns-course}

\coursemeta{Session 1 --- The neural network as an approximation class}{Session 1}{PINNs, PDEs, approximation theory, neural networks, M2 MACIA, AGM, CY Cergy Paris University}

% ================================================================
\begin{document}

\coursetitle{Session 1}{The neural network\\[3pt] as an approximation class}{}

\vspace{6pt}

\begin{keybox}
\textbf{Prerequisites.} Functional analysis (Hahn--Banach, Riesz representation), measure theory and
integration, Fourier transform of finite measures, Sobolev spaces, elementary probability.
\emph{No prior exposure to machine learning is assumed.}

\smallskip
The problem set (§\,7) is designed to be worked through independently; full solutions are in §\,8.
\end{keybox}

\bigskip

\newpage
\tableofcontents
\newpage

% ================================================================
\section{The approximation class}

\subsection{Definitions}

Throughout, $d \ge 1$ is the spatial dimension and $\sigma : \R \to \R$ is a fixed function, called
the \kw{activation function}. For a vector $y = (y_1,\dots,y_m) \in \R^m$ we write
$\sigma(y) := (\sigma(y_1),\dots,\sigma(y_m))$ for the componentwise application of $\sigma$.

\begin{definition}[Multilayer perceptron]\label{def:mlp}
Let $L \ge 1$ and $N_0 = d,\, N_1,\dots,N_L \in \N^*$. A \kw{neural network} (or multilayer
perceptron) with architecture $(N_0,\dots,N_L)$ and activation $\sigma$ is the function
$\Phi_\theta : \R^d \to \R^{N_L}$ defined by
\[
\Phi_\theta = A_L \circ \sigma \circ A_{L-1} \circ \sigma \circ \cdots \circ \sigma \circ A_1,
\qquad
A_\ell(y) = W_\ell\, y + b_\ell,\quad W_\ell \in \R^{N_\ell \times N_{\ell-1}},\ b_\ell \in \R^{N_\ell}.
\]
The parameter is $\theta = (W_1,b_1,\dots,W_L,b_L) \in \R^{P}$ with
$P = \sum_{\ell=1}^L (N_{\ell-1}N_\ell + N_\ell)$. The integer $L$ is the \kw{depth} (the number of
affine layers; there are $L-1$ hidden layers), $\max_\ell N_\ell$ is the \kw{width}, and $P$ the
number of parameters.
\end{definition}

\begin{center}
\begin{tikzpicture}[
  scale=0.95, transform shape,
  neuron/.style={circle, draw=coursgrey, minimum size=15pt, inner sep=0pt},
  inp/.style={neuron, fill=coursbg},
  hid/.style={neuron, fill=coursteal!15, draw=coursteal},
  outn/.style={neuron, fill=courswine!15, draw=courswine},
  edg/.style={-, draw=coursgrey!45, line width=0.3pt}
]
\foreach \i in {1,2}   \node[inp] (I\i) at (0, 1.2-1.2*\i) {};
\foreach \i in {1,...,4} \node[hid] (H\i) at (2.4, 2.4-1.0*\i) {};
\foreach \i in {1,...,4} \node[hid] (G\i) at (4.8, 2.4-1.0*\i) {};
\node[outn] (O1) at (7.2, -0.6) {};
\foreach \i in {1,2} \foreach \j in {1,...,4} \draw[edg] (I\i) -- (H\j);
\foreach \i in {1,...,4} \foreach \j in {1,...,4} \draw[edg] (H\i) -- (G\j);
\foreach \i in {1,...,4} \draw[edg] (G\i) -- (O1);
\node[below=2pt of I2, font=\scriptsize, text=coursgrey] {$x \in \R^2$};
\node[below=2pt of H4, font=\scriptsize, text=coursgrey] {$\sigma(A_1 x)$};
\node[below=2pt of G4, font=\scriptsize, text=coursgrey] {$\sigma(A_2 \cdot)$};
\node[below=6pt of O1, font=\scriptsize, text=coursgrey] {$\Phi_\theta(x) \in \R$};
\end{tikzpicture}
\end{center}

\begin{definition}[Network classes]
Write $\Nn^\sigma_{L,W}(d)$ for the set of functions $\R^d \to \R$ realised by a network with
activation $\sigma$, depth $\le L$ and width $\le W$, as $\theta$ ranges over $\R^P$. For $L = 2$
(one hidden layer with $n$ neurons) we simply write
\begin{equation}\label{eq:sigman}
\Sigma_n(\sigma) := \Big\{ x \mapsto \sum_{i=1}^n c_i\, \sigma(w_i \cdot x + b_i)
\ :\ c_i \in \R,\ w_i \in \R^d,\ b_i \in \R \Big\},
\qquad
\Sigma(\sigma) := \bigcup_{n \ge 1} \Sigma_n(\sigma).
\end{equation}
Thus $\Sigma(\sigma)$ is the linear span of the \kw{ridge functions}
$x \mapsto \sigma(w\cdot x + b)$.
\end{definition}

\begin{remark}
Beware of a source of confusion that pervades the literature: $\Sigma(\sigma)$, the union over $n$,
\emph{is} a vector space, whereas $\Sigma_n(\sigma)$ for \emph{fixed} $n$ is not. Density theorems
concern $\Sigma(\sigma)$; questions of cost, rate and optimisation concern $\Sigma_n(\sigma)$.
The entire difficulty of the subject lives in that gap.
\end{remark}

\subsection{Elementary structure of the class}

\begin{proposition}[Stable operations]\label{prop:stable}
Let $\sigma$ be arbitrary, $f \in \Sigma_n(\sigma)$, $g \in \Sigma_m(\sigma)$ and $\lambda \in \R$.
Then:
\begin{enumerate}[label=(\roman*)]
\item $\lambda f \in \Sigma_n(\sigma)$;
\item $f + g \in \Sigma_{n+m}(\sigma)$;
\item for every affine map $A$ from $\R^{d'}$ to $\R^d$, $f \circ A \in \Sigma_n(\sigma)$;
\item if $\sigma \in C^k(\R)$, then every $\Phi_\theta \in \Nn^\sigma_{L,W}(d)$ is of class $C^k$
and, for every multi-index $\alpha$ with $\abs{\alpha} \le k$, $\partial^\alpha \Phi_\theta$ is
computed by composition and the chain rule.
\end{enumerate}
\end{proposition}

\begin{proof}
(i) Replace $c_i$ by $\lambda c_i$. (ii) Concatenate the two sums. (iii) If $A(z) = Mz + p$ then
$\sigma(w_i\cdot(Mz+p)+b_i) = \sigma((M^{\!\top}w_i)\cdot z + (w_i\cdot p + b_i))$.
(iv) A finite composition of affine maps (which are $C^\infty$) and of $\sigma$ applied
componentwise.
\end{proof}

Item (ii) shows that $\Sigma_n$ is not stable under addition: already a hint that the class lacks
linear structure. The next result makes this precise.

\begin{proposition}[The class is not convex]\label{prop:nonconvex}
Let $d \ge 2$ and let $\sigma \in C^1(\R)$ satisfy $\sigma' > 0$ everywhere with $\sigma'$
nonconstant (for instance $\sigma = \tanh$, for which $\sigma' = \sech^2 > 0$). Then
$\Sigma_1(\sigma)$ is not convex.
\end{proposition}

\begin{proof}
Let $w_1, w_2 \in \R^d$ be linearly independent, and set
\[
f_1(x) = \sigma(w_1\cdot x), \qquad f_2(x) = \sigma(w_2 \cdot x), \qquad
f := \tfrac12 (f_1 + f_2).
\]
Then $f_1, f_2 \in \Sigma_1(\sigma)$. Suppose, for contradiction, that $f \in \Sigma_1(\sigma)$,
that is, $f(x) = c\,\sigma(w\cdot x + b)$ for some $c \in \R$, $w \in \R^d$, $b \in \R$.

First, $f$ is not constant: $\nabla f(0) = \tfrac12 \sigma'(0)(w_1 + w_2) \ne 0$, since
$\sigma'(0) > 0$ and $w_1 + w_2 \ne 0$ by linear independence. Hence $c \ne 0$ and $w \ne 0$.

For every $x$, $\nabla f(x) = c\,\sigma'(w\cdot x + b)\, w$, and $\sigma' > 0$ gives
$\nabla f(x) \in \R^*\, w$: \emph{all gradients of $f$ are collinear}. This is the characteristic
property of a ridge function: $f$ is constant on each hyperplane $\{w\cdot x = \text{const}\}$.

On the other hand, computing directly,
$\nabla f(x) = \tfrac12\big( \sigma'(w_1\cdot x)\, w_1 + \sigma'(w_2 \cdot x)\, w_2 \big)$.
By linear independence of $w_1, w_2$ there exists $x_1 \in \R^d$ with $w_1 \cdot x_1 = 0$ and
$w_2 \cdot x_1 = t$, where $t \in \R$ is chosen so that $\sigma'(t) \ne \sigma'(0)$ (possible since
$\sigma'$ is nonconstant). Then, in the basis $(w_1,w_2)$ of the plane they span,
\[
\nabla f(0) = \tfrac12 \sigma'(0)\,\big(w_1 + w_2\big),
\qquad
\nabla f(x_1) = \tfrac12 \big( \sigma'(0)\, w_1 + \sigma'(t)\, w_2\big).
\]
These two vectors are collinear if and only if
$\sigma'(0)\sigma'(t) - \sigma'(0)\sigma'(0) = 0$, that is, $\sigma'(t) = \sigma'(0)$: a
contradiction. Hence $f \notin \Sigma_1(\sigma)$, and $\Sigma_1(\sigma)$ is not convex.
\end{proof}

\begin{remark}[What is lost]\label{rem:cealost}
This seemingly innocuous proposition has far-reaching consequences. The whole Galerkin method rests
on the fact that the approximation space $V_h$ is a \emph{linear subspace}: that is what lets one
write the Euler--Lagrange equation as a linear system, invoke Céa's lemma
\[
\norm{u - u_h}_{V} \le \frac{M}{\alpha}\, \inf_{v_h \in V_h} \norm{u - v_h}_V ,
\]
and reduce error analysis to a pure approximation question. With networks none of this is available. For nonlinear fixed-size classes, such as
$\Sigma_1(\sigma)$, convexity is generally lost --- the degenerate case $L = 1$, where the
realisations are the affine maps and the class \emph{is} convex, is the exception that shows the
statement has to be qualified. The functional to be minimised is in general not convex in $\theta$
either (Exercise \ref{ex:structure}, item 4, gives a rigorous example; nonconvexity of the class
does not by itself force nonconvexity of the loss), and the minimiser may be neither unique nor
attained. This is the thread running through the course; we
shall meet it again in Session 5 (where PDE stability replaces Céa) and in Session 6 (where the gap
between ``a good $\theta$ exists'' and ``the optimiser finds it'' is quantified).
\end{remark}

\begin{remark}[Failure of closedness]\label{rem:nonclosed}
The class also fails to be closed, and this is a theorem rather than a heuristic. Fix an
architecture with $L \ge 2$ layers and at least two units in the last hidden layer. If $\sigma$ is
real analytic with bounded derivatives --- which covers $\tanh$, the logistic sigmoid and
$\arctan$ --- the set of realisations is \emph{not closed} in $W^{k,p}((0,1)^d)$ for any $k \in \N$
and any $p \in [1,\infty]$ (Mahan, King and Cloninger, 2021, Thm.\ 3.3). For $\ReLU$ the
corresponding non-closedness statement holds in $L^p$ for $1 \le p < \infty$ only: $\ReLU$ is
precisely an exception in the $L^\infty$ result of Petersen, Raslan and Voigtlaender (2021), which
is the reference for the general topological picture. The endpoint $p = \infty$ is therefore not a
detail one may quietly absorb into the quantifier.

The mechanism is the one one would guess: the difference quotients
$\lambda\big(\sigma(w\cdot x + b + 1/\lambda) - \sigma(w \cdot x + b)\big) \in \Sigma_2(\sigma)$
converge, as $\lambda \to \infty$, to $\sigma'(w\cdot x + b)$. The work in the references above is
in proving that this limit is genuinely not a realisation --- exhibiting a convergent sequence is
the easy half. Practical consequence: $\min_{\theta} J(\theta)$ may have \emph{no} solution, the
infimum being approached along diverging parameters. Weights running off to infinity are indeed
observed when training PINNs, though observing them does not by itself prove that one is in this
situation.
\end{remark}

\subsection{Regularity: a constraint the PDE imposes on the activation}

The standard activations are
\[
\ReLU(t) = \max(t,0), \qquad
\tanh(t) = \frac{\ee^t - \ee^{-t}}{\ee^t+\ee^{-t}}, \qquad
\mathrm{sig}(t) = \frac{1}{1+\ee^{-t}}, \qquad
\sin(t), \qquad
\mathrm{GELU}(t) = t\,\Phi_{\!\Nn}(t),
\]
where $\Phi_{\!\Nn}$ is the standard normal distribution function. Only the first fails to be
$C^\infty$.

\begin{proposition}[ReLU networks]\label{prop:relu}
Every function $\Phi_\theta \in \Nn^{\ReLU}_{L,W}(d)$ is continuous and piecewise affine (that is,
there is a finite partition of $\R^d$ into convex polyhedra on each of which $\Phi_\theta$ is
affine). In particular $\Phi_\theta \in W^{1,\infty}_{\mathrm{loc}}(\R^d)$ but in general
$\Phi_\theta \notin C^1$, and its distributional second derivatives are measures carried by
hyperplanes.
\end{proposition}

\begin{proof}
Being continuous and piecewise affine is preserved under composition with an affine map and under
$\ReLU$ (which cuts each polyhedron in two); induct on $L$. For the second assertion it suffices to
treat $d = 1$ and $\Phi = \ReLU$. For $\varphi \in C^\infty_c(\R)$,
\[
\scal{\Phi''}{\varphi} = \scal{\Phi}{\varphi''} = \int_0^{+\infty} x\, \varphi''(x)\,\dd x
= \big[ x\varphi'(x)\big]_0^{+\infty} - \int_0^{+\infty}\varphi'(x)\,\dd x
= 0 - \big(0 - \varphi(0)\big) = \varphi(0),
\]
so that $\ReLU'' = \delta_0$ in the sense of distributions.
\end{proof}

\begin{keybox}
\textbf{Immediate consequence for PINNs.} Consider the Poisson problem $-\Delta u = f$ and the
strong residual $\Rr(v) = -\Delta v - f$. With a ReLU activation $\Delta \Phi_\theta$ is a singular
measure, so $\Rr(\Phi_\theta)$ is not a function and its pointwise value at a collocation point is
not defined by the PDE.

\smallskip
It is worth being precise about what then happens in code, because the situation is subtler --- and
worse --- than ``the loss is meaningless''. An automatic-differentiation library does not return the
distributional derivative; it returns the classical one, which exists and equals $0$ at every point
off the finite breakpoint set. At generic collocation points the interior loss is therefore
perfectly well defined numerically, and equal to $\frac1N\sum_i \abs{f(x_i)}^2$:
\emph{a constant, with identically zero gradient in $\theta$}. Nothing crashes; the residual term
simply carries no information about $\theta$ at all. A silent failure is harder to notice than an
error message.

\smallskip
\textbf{Rule of thumb, stated carefully.} $\sigma \in C^m$ is a clean \emph{sufficient} condition
for the strong residual of an order-$m$ PDE to be defined pointwise everywhere in the classical
sense --- which is why $\tanh$ is used systematically in the PINN literature. It is not necessary in
any absolute sense: a weak formulation changes the requirement entirely. In Session 7 we shall see
that the Deep Ritz functional needs only $\Phi_\theta \in H^1$, for which a locally Lipschitz
activation on a bounded domain suffices --- $\ReLU$ qualifies, a merely continuous activation does
not. VPINNs, which integrate by parts against test functions, impose a different requirement again,
and should not be lumped together with Deep Ritz.
\end{keybox}

% ================================================================
\section{Universal approximation}

The question in this section is purely qualitative: \emph{can any continuous function be
approximated by a network, if we allow ourselves as many neurons as we wish?}

\subsection{Cybenko's theorem}

\begin{definition}
A bounded measurable function $\sigma : \R \to \R$ is \kw{sigmoidal} if
\[
\lim_{t \to +\infty} \sigma(t) = 1, \qquad \lim_{t \to -\infty} \sigma(t) = 0 .
\]
Given a compact set $K \subset \R^d$, it is \kw{discriminatory on $K$} if, for every finite
signed regular Borel measure $\mu$ on $K$,
\[
\Big( \forall w \in \R^d,\ \forall b \in \R : \int_K \sigma(w\cdot x + b)\, \dd\mu(x) = 0 \Big)
\ \Longrightarrow \ \mu = 0 .
\]
\end{definition}

\begin{lemma}[Every bounded sigmoidal function is discriminatory]\label{lem:discr}
Let $\sigma$ be bounded, measurable and sigmoidal, and let $K \subset \R^d$ be compact. Then
$\sigma$ is discriminatory on $K$.
\end{lemma}

\begin{proof}
Let $\mu$ be a finite signed Borel measure on $K$ such that
$\int_K \sigma(w\cdot x + b)\,\dd\mu(x) = 0$ for all $w, b$.

\smallskip
\noindent\emph{Step 1: $\mu$ annihilates every half-space and every hyperplane.}
Fix $w \in \R^d$, $b \in \R$, $\varphi \in \R$, and for $\lambda > 0$ set
\[
\sigma_\lambda(x) := \sigma\big(\lambda(w\cdot x + b) + \varphi\big).
\]
By hypothesis $\int_K \sigma_\lambda \,\dd\mu = 0$ for every $\lambda > 0$ (apply the hypothesis
with $\lambda w$ and $\lambda b + \varphi$). As $\lambda \to +\infty$,
\[
\sigma_\lambda(x) \longrightarrow
\begin{cases}
1 & \text{if } w\cdot x + b > 0,\\
0 & \text{if } w\cdot x + b < 0,\\
\sigma(\varphi) & \text{if } w\cdot x + b = 0.
\end{cases}
\]
The family $(\sigma_\lambda)$ is uniformly bounded by $\norm{\sigma}_\infty$ and
$\abs{\mu}(K) < \infty$; dominated convergence therefore gives, writing
$H^+_{w,b} = \{x \in K : w\cdot x + b > 0\}$ and $\Pi_{w,b} = \{x \in K : w\cdot x + b = 0\}$,
\begin{equation}\label{eq:cyb1}
\mu\big(H^+_{w,b}\big) + \sigma(\varphi)\, \mu\big(\Pi_{w,b}\big) = 0
\qquad \text{for every } \varphi \in \R .
\end{equation}
Letting $\varphi \to -\infty$ in \eqref{eq:cyb1} (so $\sigma(\varphi) \to 0$) yields
$\mu(H^+_{w,b}) = 0$; letting $\varphi \to +\infty$ yields
$\mu(H^+_{w,b}) + \mu(\Pi_{w,b}) = 0$, whence $\mu(\Pi_{w,b}) = 0$. This holds for all $w, b$.

\smallskip
\noindent\emph{Step 2: the Fourier transform of $\mu$ vanishes.}
Fix $w \in \R^d$. Since $K$ is compact, the image of $K$ under $x \mapsto w\cdot x$ is
\emph{contained in} the compact interval $[\alpha,\beta]$, where
$\alpha := \min_{x \in K} w\cdot x$ and $\beta := \max_{x \in K} w\cdot x$. (It need not \emph{be}
that interval: $K$ may be disconnected.) For bounded measurable $h : [\alpha,\beta] \to \R$ define
\[
F(h) := \int_K h(w\cdot x)\, \dd \mu(x) .
\]
$F$ is linear with $\abs{F(h)} \le \norm{h}_\infty\, \abs{\mu}(K)$. By Step 1, for every $s \in \R$,
\[
F\big(\ind_{(s,+\infty)}\big) = \mu\big(H^+_{w,-s}\big) = 0,
\qquad
F\big(\ind_{[s,+\infty)}\big) = \mu\big(H^+_{w,-s}\big) + \mu\big(\Pi_{w,-s}\big) = 0 .
\]
By subtraction and linearity, $F$ vanishes on the indicator of every interval, hence on every step
function. Every continuous function on the segment $[\alpha,\beta]$ is a uniform limit of step
functions, and $F$ is continuous for $\norm{\cdot}_\infty$; therefore $F(h) = 0$ for every
$h \in C([\alpha,\beta])$.

Applying this to $h = \cos$ and $h = \sin$:
\[
\widehat{\mu}(w) := \int_K \ee^{\,\ii\, w \cdot x}\, \dd\mu(x)
= \int_K \cos(w\cdot x)\dd\mu(x) + \ii \int_K \sin(w\cdot x)\dd\mu(x) = 0
\qquad \text{for every } w \in \R^d .
\]
The Fourier--Stieltjes transform of a finite signed measure determines that measure; hence
$\mu = 0$.
\end{proof}

\begin{theorem}[Cybenko, 1989]\label{th:cybenko}
Let $\sigma \in C(\R)$ be bounded and sigmoidal, and let $K \subset \R^d$ be compact. Then
$\Sigma(\sigma)$ is dense in $\big(C(K), \norm{\cdot}_\infty\big)$: for every $f \in C(K)$ and every
$\varepsilon > 0$ there exist $n \in \N^*$ and $(c_i, w_i, b_i)_{1 \le i \le n}$ such that
\[
\sup_{x \in K}\ \Big| f(x) - \sum_{i=1}^n c_i \,\sigma(w_i \cdot x + b_i)\Big| < \varepsilon .
\]
\end{theorem}

\begin{proof}
$S := \Sigma(\sigma)$ is a linear subspace of $C(K)$. Suppose, for contradiction, that its closure
$\overline{S}$ is a \emph{proper} closed subspace of $C(K)$.

By the Hahn--Banach theorem (geometric form, or the separation corollary) there is a continuous
linear functional $L \in C(K)'$, $L \ne 0$, with $L|_{\overline{S}} = 0$.

By the Riesz--Markov representation theorem there is a finite signed regular Borel measure $\mu$ on
$K$ such that
\[
L(f) = \int_K f \,\dd \mu \qquad \text{for every } f \in C(K),
\qquad \text{with } \norm{L}_{C(K)'} = \abs{\mu}(K) .
\]
Since $L \ne 0$, we have $\mu \ne 0$.

But for all $w \in \R^d$ and $b \in \R$ the function $x \mapsto \sigma(w\cdot x + b)$ belongs to
$S$, hence
\[
\int_K \sigma(w\cdot x + b)\,\dd\mu(x) = L\big(\sigma(w\cdot\, \cdot + b)\big) = 0 .
\]
Lemma \ref{lem:discr} then gives $\mu = 0$: a contradiction. Therefore $\overline{S} = C(K)$.
\end{proof}

\begin{remark}[A step not to skip]
No reduction to a cube is needed above, because Lemma \ref{lem:discr} was stated for an arbitrary
compact $K$ and its proof uses only compactness. Had we stated the lemma on $[0,1]^d$ alone --- as
Cybenko does --- the passage to a general compact $K \subset [0,1]^d$ would require extending $\mu$
from $K$ to the cube by zero. That extension is harmless, but it is a step, and it is routinely
omitted in secondary accounts.
\end{remark}

\begin{remark}[What the theorem does not say]\label{rem:cybenko-limits}
Three warnings, worth repeating to students:
\begin{enumerate}[label=(\alph*)]
\item The proof is \emph{nonconstructive}: it argues by contradiction via Hahn--Banach. It produces
no network and no algorithm.
\item It provides \emph{no rate}: nothing is said about how many neurons are needed to reach
$\varepsilon$. That is the subject of §\,3.
\item It says nothing about whether such $(c_i,w_i,b_i)$ can be \emph{found} by an optimisation
method. That is the subject of Session 6, and it is precisely where the real gap between the theory
and the practice of PINNs lies.
\end{enumerate}
\end{remark}

\subsection{The complete characterisation: the Leshno--Lin--Pinkus--Schocken theorem}

Cybenko's theorem assumes $\sigma$ sigmoidal, which excludes $\ReLU$ (unbounded). The optimal
characterisation is due to Leshno, Lin, Pinkus and Schocken.

\begin{theorem}[LLPS, 1993]\label{th:llps}
Let $\sigma \in C(\R)$. Then $\Sigma(\sigma)$ is dense in $C(K)$ for every compact
$K \subset \R^d$ and every $d \ge 1$ \emph{if and only if} $\sigma$ is not a polynomial.
\end{theorem}

We prove the necessary condition in full, and then the sufficient condition in the case
$\sigma \in C^\infty$ --- the general $C^0$ case reduces to it by mollification, see Pinkus (1999,
§\,3).

\begin{proposition}[Necessity]\label{prop:poly}
If $\sigma$ is a polynomial of degree $m$, then $\Sigma(\sigma)$ is dense in $C(K)$ for no compact
set $K$ with nonempty interior.
\end{proposition}

\begin{proof}
If $\sigma(t) = \sum_{j=0}^m a_j t^j$, then for all $w,b$ the function
$x \mapsto \sigma(w\cdot x + b) = \sum_j a_j (w\cdot x + b)^j$ is a polynomial in $x$ of total
degree $\le m$. Hence $\Sigma(\sigma) \subseteq \Pp_m(\R^d)$, the space of polynomials of total
degree $\le m$. But $\Pp_m(\R^d)$ is finite-dimensional, hence closed in $C(K)$. Therefore
$\overline{\Sigma(\sigma)} \subseteq \Pp_m(\R^d)$, which is a proper subspace of $C(K)$ as soon as
$K$ has nonempty interior (the function $x \mapsto \abs{x_1 - x_1^0}$, for $x^0$ interior, is not
polynomial near $x^0$). Hence $\Sigma(\sigma)$ is not dense.
\end{proof}

\begin{proposition}[Sufficiency, smooth case]\label{prop:llps-smooth}
Let $\sigma \in C^\infty(\R)$ be nonpolynomial. Then $\Sigma(\sigma)$ is dense in $C(K)$ for every
compact $K \subset \R^d$.
\end{proposition}

\begin{proof}
\emph{Step 1: ridge monomials lie in the closure.}
Write $V := \overline{\Sigma(\sigma)}$ (closure in $C(K)$), a closed linear subspace. Fix
$w \in \R^d$ and $b \in \R$. For $\lambda \in \R$ and $h \ne 0$,
\[
x \longmapsto \frac{\sigma\big((\lambda + h)\, w\cdot x + b\big) - \sigma\big(\lambda\, w \cdot x + b\big)}{h}
\ \in\ \Sigma_2(\sigma) \subseteq V .
\]
As $h \to 0$, this expression converges \emph{uniformly on $K$} (since $w\cdot x$ is bounded on $K$
and $\sigma'$ is uniformly continuous on compact sets) to
$x \mapsto (w\cdot x)\, \sigma'(\lambda\, w\cdot x + b)$. Since $V$ is closed, this function lies in
$V$.

Iterating the same argument $k$ times (differentiating in $\lambda$, which is legitimate throughout
by uniform convergence on the compact set $K$) gives
\[
x \longmapsto (w\cdot x)^k\, \sigma^{(k)}\big(\lambda\, w\cdot x + b\big) \ \in\ V
\qquad \text{for all } k \in \N,\ \lambda \in \R,\ b \in \R .
\]
Taking $\lambda = 0$: $x \mapsto \sigma^{(k)}(b)\, (w\cdot x)^k \in V$.

Since $\sigma$ is nonpolynomial and $\sigma \in C^\infty$, for every $k$ there is $b_k \in \R$ with
$\sigma^{(k)}(b_k) \ne 0$ --- otherwise $\sigma^{(k)} \equiv 0$ for some $k$ and $\sigma$ would be
polynomial. Hence
\[
x \longmapsto (w\cdot x)^k \ \in\ V \qquad \text{for all } k \in \N,\ w \in \R^d .
\]

\smallskip
\noindent\emph{Step 2: ridge monomials span all polynomials.}
Fix $k$ and $i_1,\dots,i_k \in \{1,\dots,d\}$. The polarisation identity
\[
k!\, y_1 y_2 \cdots y_k \;=\; \sum_{S \subseteq \{1,\dots,k\}} (-1)^{k - \abs{S}}
\Big( \sum_{j \in S} y_j \Big)^{k}
\]
(valid in any commutative ring; expand by the multinomial formula and observe that only the terms
containing each $y_j$ exactly once survive) applied to $y_j = x_{i_j}$ gives
\[
k!\; x_{i_1}\cdots x_{i_k} \;=\; \sum_{S} (-1)^{k-\abs{S}} \big( w_S \cdot x \big)^k,
\qquad w_S := \sum_{j \in S} e_{i_j}.
\]
Each term on the right lies in $V$ by Step 1, hence $x_{i_1}\cdots x_{i_k} \in V$. By linearity $V$
contains every monomial, hence every polynomial.

\smallskip
\noindent\emph{Step 3: conclusion.}
$\Pp(\R^d)$ is a subalgebra of $C(K)$ containing the constants and separating points; by the
Stone--Weierstrass theorem it is dense in $C(K)$. Since $\Pp(\R^d) \subseteq V$ and $V$ is closed,
$V = C(K)$.
\end{proof}

\begin{remark}
Proposition \ref{prop:llps-smooth} is also the key to a useful exercise (Exercise \ref{ex:poly}):
it shows that the real obstruction to universality is not the shape of $\sigma$ but its
\emph{degree of nonpolynomiality}. A network with activation $\sigma(t) = t^2$ and depth $L$ spans
only polynomials of degree $\le 2^{L-1}$: depth raises the degree, but one stays inside a
polynomial class.
\end{remark}

% ================================================================
\section{Quantitative bounds: the Barron class}

Cybenko's theorem gives no rate. The quantitative question is: how many neurons are needed to reach
accuracy $\varepsilon$? For a \emph{linear} approximation method (finite elements, polynomials,
wavelets) on a ball of $W^{s,\infty}([0,1]^d)$, the answer is given by the Kolmogorov $n$-widths:
\begin{equation}\label{eq:kolm}
d_n\big( B_{W^{s,\infty}}, L^\infty \big) \;\asymp\; n^{-s/d},
\end{equation}
that is, $n \asymp \varepsilon^{-d/s}$: the \kw{curse of dimensionality}. The next result shows
that a well-chosen class of functions escapes this bound \emph{with networks}, at the price of a
strong structural hypothesis.

\subsection{Maurey's lemma}

This is a purely Hilbertian result, independent of networks, and remarkably simple.

\begin{lemma}[Maurey; Pisier; Jones; Barron]\label{lem:maurey}
Let $\Hh$ be a Hilbert space and $G \subseteq \Hh$ with
$\displaystyle\sup_{g \in G} \norm{g}_\Hh \le \beta < \infty$. Let $f \in \co(G)$, the convex hull
of $G$. Then for every $n \ge 1$ there exist $g_1,\dots,g_n \in G$ such that
\[
\Big\| f - \frac1n \sum_{i=1}^n g_i \Big\|_{\Hh}^2 \;\le\; \frac{\beta^2 - \norm{f}_\Hh^2}{n} .
\]
If $f \in \overline{\co}(G)$, then for every $\varepsilon > 0$ there exist $g_1,\dots,g_n \in G$
with $\norm{f - \frac1n\sum_i g_i}_\Hh \le \beta/\sqrt{n} + \varepsilon$.
\end{lemma}

\begin{proof}
Write $f = \sum_{j=1}^m \lambda_j h_j$ with $h_j \in G$, $\lambda_j \ge 0$, $\sum_j \lambda_j = 1$.
Let $G_1, \dots, G_n$ be independent identically distributed random variables with values in
$\{h_1,\dots,h_m\}$ and $\Prob(G_i = h_j) = \lambda_j$. Then $\E[G_i] = f$ (a Bochner expectation,
here a finite sum).

Set $S_n = \frac1n \sum_{i=1}^n G_i$. The variables $G_i - f$ are independent, centred and
$\Hh$-valued, so the cross terms vanish:
\[
\E \norm{f - S_n}^2
= \frac{1}{n^2}\, \E\Big\| \sum_{i=1}^n (G_i - f) \Big\|^2
= \frac{1}{n^2} \sum_{i=1}^n \E\norm{G_i - f}^2
= \frac{1}{n}\, \E\norm{G_1 - f}^2 .
\]
Moreover $\E\norm{G_1 - f}^2 = \E\norm{G_1}^2 - 2\scal{\E G_1}{f} + \norm{f}^2
= \E\norm{G_1}^2 - \norm{f}^2 \le \beta^2 - \norm{f}^2$. Hence
$\E\norm{f - S_n}^2 \le (\beta^2 - \norm{f}^2)/n$.

A random variable is at most its expectation with positive probability: there is therefore a
realisation $(g_1,\dots,g_n)$ for which
$\norm{f - \frac1n\sum_i g_i}^2 \le (\beta^2 - \norm f^2)/n$.

For $f \in \overline{\co}(G)$, choose $\tilde f \in \co(G)$ with
$\norm{f - \tilde f} \le \varepsilon$, apply the above to $\tilde f$, and use the triangle
inequality.
\end{proof}

\begin{remark}
The rate $n^{-1/2}$ is \emph{independent of the dimension}: this is the whole novelty. It is not, in
fact, a phenomenon specific to neural networks but a general property of \emph{nonlinear} $n$-term
approximation in a Hilbert space. The price to pay is that the membership $f \in \overline{\co}(G)$
is a strong hypothesis: it is exactly what defines the Barron class.
\end{remark}

\subsection{The Barron class}

\begin{definition}[Barron spectral seminorm]\label{def:barron}
Let $r > 0$ and $B_r = \{x \in \R^d : \abs{x} \le r\}$. For $f : B_r \to \R$ set
\[
C_f := \inf \Big\{ \int_{\R^d} \abs{\xi}\,\dd\abs{\nu}(\xi) \Big\},
\]
the infimum being over all finite complex Borel measures $\nu$ on $\R^d$ such that
\[
f(x) = \int_{\R^d} \ee^{\,\ii\, \xi\cdot x}\, \dd\nu(\xi)
\qquad \text{for every } x \in B_r ,
\]
with $C_f := +\infty$ if no such $\nu$ exists. Write $\Gamma_C := \{ f : C_f \le C\}$ for the
\kw{Barron class} with constant $C$.
\end{definition}

\begin{remark}[Three points of hygiene]\label{rem:barron-hygiene}
\emph{(a) Measures, not densities.} Requiring $\dd\nu = \tilde f\,\dd\xi$ with
$\tilde f \in L^1$ would be a real restriction, not a harmless convenience: a pure frequency such
as $x \mapsto \sin(K x_1)$ has no $L^1$ Fourier density, and mollifying an atom destroys the
\emph{exact} identity on $B_r$ that the definition demands. With measures, $\nu$ may have atoms and
the examples below are legitimate as they stand. This is also how Barron states it.

\emph{(b) It is a seminorm.} Every constant function has $C_f = 0$: take $\nu = c\,\delta_0$, for
which $\int\abs\xi\,\dd\abs\nu = 0$. This is exactly why $f(0)$ is subtracted in
Theorem \ref{th:barron}: the constant part of $f$ is invisible to $C_f$ and must be supplied
separately.

\emph{(c) Exhibiting a representation gives only an upper bound.} Since the identity is required
only on $B_r$, a particular $\nu$ yields $C_f \le \int\abs\xi\,\dd\abs\nu$ and nothing more: a
cheaper representation agreeing with $f$ on $B_r$ and differing elsewhere may well exist. Any
estimate obtained solely by exhibiting a representing measure is therefore an upper bound, unless a
matching lower bound is proved --- as it is in Proposition \ref{prop:barron-elem}\,(iv), where
part (ii) supplies it. Upper bounds are in any case all Theorem \ref{th:barron} needs, since the
bound it proves is increasing in $C_f$. For the same reason the infimum need not be attained, and
the proof below carries an $\eta > 0$ throughout.
\end{remark}

\begin{proposition}[Elementary properties]\label{prop:barron-elem}
\begin{enumerate}[label=(\roman*)]
\item $\Gamma_C$ is convex and symmetric, and $C_{\lambda f} = \abs{\lambda} C_f$.
\item If $C_f < \infty$ then $f$ is Lipschitz on $B_r$ with constant $\le C_f$.
\item For $f(x) = \ee^{-\abs{x}^2/2}$ on $\R^d$,
$C_f \le \E\abs{Z} = \sqrt2\,\Gamma\!\big(\tfrac{d+1}{2}\big)/\Gamma\!\big(\tfrac d2\big) \sim \sqrt d$,
where $Z \sim \Nn(0,I_d)$.
\item For $f_K(x) = \sin(K x_1)$ one has $C_{f_K} = K$ --- the one entry in this list where the
seminorm can be computed exactly, the lower bound coming from (ii). So the seminorm does grow
linearly with frequency, and this is proved, not merely bounded above.
\end{enumerate}
\end{proposition}

\begin{proof}
(i) Immediate from the definition. (ii) For any representing measure $\nu$ and any $x,y \in B_r$, using
$\abs{\ee^{\ii a} - \ee^{\ii b}} \le \abs{a-b}$,
\[
\abs{f(x)-f(y)} = \Big|\int \big(\ee^{\ii \xi\cdot x} - \ee^{\ii \xi\cdot y}\big)\dd\nu(\xi)\Big|
\le \int \abs{\xi\cdot(x-y)}\,\dd\abs{\nu}(\xi)
\le \abs{x-y} \int\abs{\xi}\dd\abs{\nu}(\xi) ;
\]
taking the infimum over $\nu$ gives the Lipschitz constant $C_f$.
(iii) Take $\dd\nu(\xi) = (2\pi)^{-d/2}\ee^{-\abs{\xi}^2/2}\dd\xi$, which represents $f$ on all
of $\R^d$ since
$\int_{\R^d} \ee^{\ii \xi x}(2\pi)^{-d/2}\ee^{-\abs\xi^2/2}\dd\xi = \ee^{-\abs x^2/2}$. Then
\[
C_f \le \int \abs\xi\, (2\pi)^{-d/2} \ee^{-\abs\xi^2/2}\,\dd\xi = \E\abs{Z},
\]
and the formula for $\E\abs Z$ is that of the $\chi$ distribution with $d$ degrees of freedom; the
equivalent $\sqrt d$ follows from Stirling's formula.
(iv) For the upper bound, take $\nu = \frac{1}{2\ii}\big(\delta_{Ke_1} - \delta_{-Ke_1}\big)$, a
finite complex measure with
$\int \ee^{\ii\xi\cdot x}\dd\nu(\xi) = \frac{1}{2\ii}\big(\ee^{\ii Kx_1} - \ee^{-\ii Kx_1}\big)
= \sin(Kx_1)$. Its total variation is
$\abs\nu = \tfrac12\big(\delta_{Ke_1} + \delta_{-Ke_1}\big)$, so
$C_{f_K} \le \int\abs\xi\dd\abs\nu = \tfrac12 K + \tfrac12 K = K$.

For the lower bound, apply (ii): $C_{f_K}$ dominates the Lipschitz constant of $f_K$ on $B_r$, and
$\nabla f_K(x) = K\cos(Kx_1)\,e_1$, so $\abs{\nabla f_K(0)} = K$ and the Lipschitz constant is at
least $K$. Hence $C_{f_K} \ge K$, and therefore $C_{f_K} = K$.
\end{proof}

\begin{remark}[Finite seminorm does impose regularity]\label{rem:barron-c1}
It is tempting to read (iv) as saying that the Barron class ignores regularity altogether. It does
not. If $\int\abs\xi\dd\abs\nu < \infty$ then one may differentiate under the integral sign,
\[
\partial_j f(x) = \int_{\R^d} \ii\, \xi_j\, \ee^{\ii \xi\cdot x}\, \dd\nu(\xi),
\]
the domination being by $\abs{\xi} \in L^1(\abs\nu)$; the same domination makes each $\partial_j f$
continuous. So a function of finite spectral Barron seminorm admits a $C^1$ extension to $\R^d$. The
correct statement is therefore not that the class is indifferent to smoothness, but that it is not
\emph{graded} by classical smoothness order: an analytic function can have a large seminorm, and a
finite seminorm already forces $C^1$.
\end{remark}

\begin{remark}\label{rem:barron-freq}
Item (iv) is worth remembering for Session 6: in the bounds of this section, and again in the
analysis of training, it is the \emph{frequency} content of the solution rather than its smoothness
that drives the estimate. Stated carefully: a high frequency makes the Barron bound expensive --- an
upper bound, so this is not a lower bound on the number of neurons --- and is separately associated
with slower training in the regimes where that can be analysed. A spectral method may handle the
same function comfortably, but only if its basis is adapted to it; a high frequency still demands
adequate resolution there too.
\end{remark}

\subsection{Barron's theorem}

\begin{theorem}[Barron, 1993]\label{th:barron}
Let $\sigma$ be sigmoidal with $0 \le \sigma \le 1$, let $r > 0$, and let $\mu$ be a probability
measure on $B_r$. Let $f : B_r \to \R$ with $C_f < \infty$. Then, for every $n \ge 1$ and every
$\varepsilon > 0$, there exist $c_1,\dots,c_n \in \R$, $w_i \in \R^d$, $b_i \in \R$ with
$\sum_{i=1}^n \abs{c_i} \le 2 r\, C_f + \varepsilon$, such that
\[
\Big\| f - f(0) - \sum_{i=1}^n c_i\, \sigma(w_i\cdot x + b_i) \Big\|_{L^2(\mu)}
\;\le\; \frac{2\, r\, C_f}{\sqrt n} \;+\; \varepsilon .
\]
\end{theorem}

The proof has two stages: first one shows that $f - f(0)$ lies in the closed convex hull of a family
of uniformly bounded \emph{ridge sinusoids}, then that each of these sinusoids itself lies in the
closed convex hull of the sigmoidal units. Maurey's lemma concludes.

\begin{lemma}[Stage A: a mixture of ridge sinusoids]\label{lem:barron-A}
Let $f$ be as in Theorem \ref{th:barron} and let $\eta > 0$. Unless $f$ is constant on $B_r$, there
are a number $\Lambda_\eta \in (0,\, C_f + \eta]$, a probability measure $P$ on
$\R^d \setminus \{0\}$, and a family $\big(g(\cdot,\xi)\big)_\xi$ of functions on $B_r$ such that
\[
f(x) - f(0) = \int_{\R^d\setminus\{0\}} g(x,\xi)\, \dd P(\xi)
\quad \text{for every } x \in B_r,
\qquad
\sup_{\xi}\ \norm{g(\cdot,\xi)}_{L^\infty(B_r)} \le r\, \Lambda_\eta,
\]
and each $g(\cdot,\xi)$ has the form $x \mapsto h_\xi(a_\xi \cdot x)$ with $a_\xi \in \R^d$ a unit
vector and $h_\xi : [-r,r]\to\R$ Lipschitz with constant $\le \Lambda_\eta$ and $h_\xi(0)=0$.
\end{lemma}

\begin{proof}
The infimum in Definition \ref{def:barron} need not be attained, so we fix $\eta > 0$ and choose a
representing measure $\nu$ with $\int\abs\xi\,\dd\abs\nu \le C_f + \eta$. Let
$\dd\nu = \ee^{\ii\theta(\xi)}\dd\abs\nu(\xi)$ be its polar decomposition. Since $f$ is real
valued,
\[
f(x) - f(0)
= \operatorname{Re} \int \big(\ee^{\ii \xi\cdot x} - 1\big)\dd\nu(\xi)
= \int \Big( \cos\big(\xi\cdot x + \theta(\xi)\big) - \cos\theta(\xi) \Big)\, \dd\abs{\nu}(\xi) .
\]
Set $\Lambda_\eta := \int_{\xi \ne 0} \abs\xi\, \dd\abs{\nu}(\xi) \le C_f + \eta$; note that any
atom of $\nu$ at $\xi = 0$ contributes to neither side, the integrand vanishing there. If
$\Lambda_\eta = 0$ then $\abs\nu$ is carried by $\{0\}$ and $f$ is constant on $B_r$, the case
excluded by hypothesis. Otherwise put
\[
\dd P(\xi) := \frac{\abs\xi\,\dd\abs{\nu}(\xi)}{\Lambda_\eta}
\quad\text{(a probability measure on } \R^d\setminus\{0\}),
\qquad
g(x,\xi) := \frac{\Lambda_\eta}{\abs\xi}\Big( \cos\big(\xi\cdot x + \theta(\xi)\big) - \cos\theta(\xi)\Big).
\]
Then $f(x) - f(0) = \int g(x,\xi)\dd P(\xi)$ by construction. Moreover, since
$\abs{\cos(u+v) - \cos u} \le \abs v$,
\[
\abs{g(x,\xi)} \le \frac{\Lambda_\eta}{\abs\xi}\,\abs{\xi\cdot x} \le \Lambda_\eta\, \abs x \le r\, \Lambda_\eta
\qquad \text{for } x \in B_r .
\]
Finally, setting $a_\xi := \xi/\abs\xi$ and
$h_\xi(z) := \frac{\Lambda_\eta}{\abs\xi}\big(\cos(\abs\xi z + \theta(\xi)) - \cos\theta(\xi)\big)$,
we have $g(x,\xi) = h_\xi(a_\xi\cdot x)$, $h_\xi(0) = 0$, and
$\abs{h_\xi'(z)} = \frac{\Lambda_\eta}{\abs\xi}\,\abs\xi\, \abs{\sin(\abs\xi z + \theta)} \le \Lambda_\eta$.
\end{proof}

\begin{lemma}[Stage B: a Lipschitz function is a mixture of steps]\label{lem:barron-B}
Let $h : [-r,r] \to \R$ be Lipschitz with constant $\le \Lambda$ and $h(0) = 0$. Then, for every
$z \in [-r,r]$,
\[
h(z) = \int_{-r}^{r} h'(u)\, s_u(z)\, \dd u,
\qquad \text{where } \
s_u(z) := \begin{cases} \ind_{\{z > u\}} & \text{if } u \ge 0,\\[2pt]
-\,\ind_{\{z < u\}} & \text{if } u < 0. \end{cases}
\]
In particular $h$ belongs to
$\overline{\co}\big( \{\pm 2r\Lambda\, s_u : u \in [-r,r]\} \cup \{0\}\big)$ for convergence in
$L^2$ of any finite measure on $[-r,r]$.
\end{lemma}

\begin{proof}
$h$ is Lipschitz, hence absolutely continuous, and $h(z) = \int_0^z h'(u)\dd u$ with
$\abs{h'} \le \Lambda$ almost everywhere. If $z \ge 0$:
\[
\int_{-r}^r h'(u)s_u(z)\dd u
= \int_{0}^{r} h'(u)\ind_{\{z>u\}}\dd u - \int_{-r}^{0} h'(u)\ind_{\{z<u\}}\dd u
= \int_0^z h'(u)\dd u - 0 = h(z),
\]
the second integral vanishing because $u < 0 \le z$ rules out $z < u$. If $z < 0$:
\[
\int_{-r}^r h'(u)s_u(z)\dd u = 0 - \int_{z}^{0} h'(u)\dd u = h(z) - h(0) = h(z).
\]
For the second assertion: $\int_{-r}^r \abs{h'(u)}\dd u \le 2r\Lambda$. Setting
$m := \int \abs{h'}$ and, if $m>0$, $\dd Q(u) := \abs{h'(u)}\dd u / m$, we get
\[
h(z) = m \int \operatorname{sgn}\big(h'(u)\big)\, s_u(z)\, \dd Q(u)
= \int \Big[ \tfrac{m}{2r\Lambda}\cdot \operatorname{sgn}(h'(u))\, 2r\Lambda\, s_u(z)
+ \big(1 - \tfrac{m}{2r\Lambda}\big)\cdot 0 \Big]\dd Q(u),
\]
which is a mixture (an integral against a probability measure) of elements of the stated family.
Such a mixture belongs to the closed convex hull: this is exactly the content of the variance
computation in Lemma \ref{lem:maurey}, applied to $n$ i.i.d.\ draws from $Q$.
\end{proof}

\begin{proof}[Proof of Theorem \ref{th:barron}]
We work in the Hilbert space $\Hh = L^2(\mu)$, where $\mu$ is a probability measure on $B_r$. If
$f$ is constant on $B_r$ the statement is trivial, so assume otherwise, fix $\eta > 0$ and let
$\Lambda_\eta \le C_f + \eta$ be as in Lemma \ref{lem:barron-A}. Set
\[
G_\sigma := \big\{ x \mapsto \pm\, 2r \Lambda_\eta\, \sigma(w\cdot x + b)\ :\ w \in \R^d,\ b \in \R \big\} \cup \{0\}.
\]
Since $0 \le \sigma \le 1$ and $\mu$ is a probability measure,
$\sup_{g \in G_\sigma}\norm{g}_{L^2(\mu)} \le 2r\Lambda_\eta =: \beta$.

\smallskip\noindent
\emph{(1) Steps lie in $\overline{\co}(G_\sigma)$.}
Let $a \in \R^d$ be a unit vector and $u \in [-r,r]$. As $\lambda \to +\infty$,
\[
\sigma\big(\lambda(a\cdot x - u)\big) \longrightarrow \ind_{\{a\cdot x > u\}},
\qquad
-\,\sigma\big(\lambda(u - a\cdot x)\big) \longrightarrow -\,\ind_{\{a\cdot x < u\}},
\]
pointwise at every $x$ with $a\cdot x \ne u$. Both approximants are of the form
$\pm\sigma(w\cdot x + b)$, so their multiples by $2r\Lambda_\eta$ lie in $G_\sigma$; in particular
\emph{no additive constant is needed}, which is why the two cases are treated separately rather
than through $\ind_{\{a\cdot x<u\}} = 1 - \ind_{\{a\cdot x>u\}}$. If $\mu(\{x : a\cdot x = u\}) = 0$,
dominated convergence upgrades the pointwise convergence to convergence in $L^2(\mu)$, so
$\pm 2r\Lambda_\eta\, s_u \in \overline{G_\sigma}$. The set of $u \in [-r,r]$ for which
$\mu(\{a\cdot x = u\}) > 0$ is at most countable, hence Lebesgue-null: it does not affect the
integral in Lemma \ref{lem:barron-B}.

\smallskip\noindent
\emph{(2) Each $g(\cdot,\xi)$ lies in $\overline{\co}(G_\sigma)$.}
By Lemma \ref{lem:barron-A}, $g(\cdot,\xi) = h_\xi(a_\xi \cdot x)$ with $h_\xi$ Lipschitz of
constant $\le \Lambda_\eta$ and $h_\xi(0)=0$; by Lemma \ref{lem:barron-B} (with
$\Lambda = \Lambda_\eta$), $h_\xi$ is a mixture of functions $\pm 2r\Lambda_\eta\, s_u$, hence, by
(1) together with convexity and closedness of $\overline{\co}(G_\sigma)$,
$g(\cdot,\xi) \in \overline{\co}(G_\sigma)$.

\smallskip\noindent
\emph{(3) $f - f(0) \in \overline{\co}(G_\sigma)$.}
By Lemma \ref{lem:barron-A}, $f - f(0) = \int g(\cdot,\xi)\dd P(\xi)$, a mixture of elements of the
closed convex set $\overline{\co}(G_\sigma)$; by the same variance argument as in
Lemma \ref{lem:maurey}, this mixture is a limit in $L^2(\mu)$ of finite averages of elements of
$\overline{\co}(G_\sigma)$, hence belongs to it.

\smallskip\noindent
\emph{(4) Conclusion.} Maurey's lemma \ref{lem:maurey} applied with $\Hh = L^2(\mu)$,
$G = G_\sigma$, $\beta = 2r\Lambda_\eta$ and $f - f(0) \in \overline{\co}(G_\sigma)$ yields
$g_1,\dots,g_n \in G_\sigma$ with
\[
\Big\| f - f(0) - \frac1n \sum_{i=1}^n g_i \Big\|_{L^2(\mu)}
\le \frac{2r\Lambda_\eta}{\sqrt n} + \frac{\varepsilon}{2}
\le \frac{2rC_f}{\sqrt n} + \frac{2r\eta}{\sqrt n} + \frac{\varepsilon}{2} .
\]
Choosing $\eta \le \varepsilon/(4r)$ at the outset absorbs the middle term, giving the stated
bound. Each $g_i$ is of the form $\pm 2r\Lambda_\eta\,\sigma(w_i\cdot x + b_i)$ (or $0$); setting
$c_i = \pm 2r\Lambda_\eta/n$ gives the stated form, with
$\sum_i \abs{c_i} \le 2r\Lambda_\eta \le 2rC_f + \varepsilon$.
\end{proof}

\begin{keybox}
\textbf{Reading the theorem.} To reach $L^2$ error $\varepsilon$ on $\Gamma_C$ it suffices to take
$n = O\big((rC/\varepsilon)^2\big)$ neurons, \emph{whatever the dimension $d$}. The theorem provides
a sufficient width; it asserts neither an equivalence nor a lower bound. Compare with
\eqref{eq:kolm}, where the two-sided $n \asymp \varepsilon^{-d/s}$ \emph{is} legitimate, being read
off a Kolmogorov $n$-width --- a quantity defined as an infimum over all linear methods.

\smallskip
\textbf{Two essential caveats.} \emph{(a)} The constant $C_f$ may itself grow with $d$; it grows
like $\sqrt d$ for the Gaussian (Proposition \ref{prop:barron-elem}), but can grow exponentially
for other families. The absence of the curse is therefore \emph{conditional on a structural
hypothesis on $f$}, not free. \emph{(b)} The theorem is an \emph{existence} result: it says nothing
about how to find the $(c_i,w_i,b_i)$. We shall see in Session 6 that gradient descent does not, in
general, find them.
\end{keybox}

% ================================================================
\section{Approximation in Sobolev norms}

\subsection{Why this is the right question for PINNs}

A PINN does not minimise $\norm{u - \Phi_\theta}$ but $\norm{\Dd[\Phi_\theta] - f}$. For an
operator of order $m$ one must therefore control $\Phi_\theta$ \emph{and its derivatives up to
order $m$}: what we need are approximation statements in $W^{m,p}$, not in $L^p$ or $C^0$. This is
a strictly harder question: approximating a function in the uniform norm implies nothing about its
derivatives.

\subsection{The mechanism: emulating monomials by finite differences}

The following elementary lemma contains the idea behind many of the modern constructions.

\begin{lemma}[Emulating $x^2$]\label{lem:emul}
Let $\sigma \in C^4(\R)$ and let $a \in \R$ satisfy $\sigma''(a) \ne 0$. For $h > 0$ set
\[
\Psi_h(x) := \frac{\sigma(a + hx) + \sigma(a - hx) - 2\sigma(a)}{h^2\, \sigma''(a)} .
\]
Then $\Psi_h$ is realised by a network with one hidden layer of three neurons, and for every
$R > 0$ there is $\kappa = \kappa(R,\sigma,a)$ such that, for $0 < h \le 1$,
\[
\sup_{\abs x \le R} \abs{\Psi_h(x) - x^2} \le \kappa\, h^2,
\qquad
\sup_{\abs x \le R} \abs{\Psi_h'(x) - 2x} \le \kappa\, h^2 .
\]
\end{lemma}

\begin{proof}
Taylor's formula with Lagrange remainder at order 4 around $a$ gives, for $\abs{u} \le hR$,
\[
\sigma(a+u) + \sigma(a-u) - 2\sigma(a) = \sigma''(a)\, u^2 + \frac{u^4}{24}\,\big(\sigma^{(4)}(\eta_+) + \sigma^{(4)}(\eta_-)\big)
\]
for some $\eta_\pm$ between $a$ and $a \pm u$ (the odd-order terms cancel; the coefficient is
$1/4! = 1/24$, once for each of the two expansions). With $u = hx$ and
$M_4 := \sup_{\abs{t - a} \le R}\abs{\sigma^{(4)}(t)}$:
\[
\abs{\Psi_h(x) - x^2} \le \frac{h^4 x^4}{24\, h^2 \abs{\sigma''(a)}} \cdot 2 M_4
\le \frac{R^4 M_4}{12\,\abs{\sigma''(a)}}\; h^2 .
\]
For the derivative, $\Psi_h'(x) = \dfrac{\sigma'(a+hx) - \sigma'(a-hx)}{h\,\sigma''(a)}$, and
Taylor at order 3 applied to $\sigma'$ gives
\[
\sigma'(a+u)-\sigma'(a-u) = 2u\,\sigma''(a) + \frac{u^3}{6}\big(\sigma^{(4)}(\zeta_+) + \sigma^{(4)}(\zeta_-)\big),
\]
whence $\abs{\Psi_h'(x) - 2x} \le \dfrac{R^3 M_4}{3\abs{\sigma''(a)}}\, h^2$. Take $\kappa$ to be the
maximum of the two constants.
\end{proof}

\begin{remark}[The price: enormous weights]\label{rem:weights}
The network $\Psi_h$ has output weights of size $1/(h^2\abs{\sigma''(a)})$: to gain a factor $h^2$
in accuracy one pays a growth of $h^{-2}$ in the parameters. This is not an artefact of the
construction: many constructive approximation results in Sobolev norm by smooth networks exhibit
some form of this trade-off. Weight size is not
a detail --- it appears explicitly in the generalisation bounds of Session 5, where it controls the
Lipschitz constant of the residual, hence the quadrature error.
\end{remark}

\begin{remark}
For $\sigma = \tanh$ one has $\sigma''(t) = -2\tanh(t)\sech^2(t)$, which vanishes at $0$: one must
therefore take $a \ne 0$, for instance $a = 1$ where $\sigma''(1) \approx -0.64$. The fact that
$\tanh$ is odd explains why the constructions in the literature work around an off-centre point.
\end{remark}

\subsection{The reference statements}

\begin{theorem}[De Ryck, Lanthaler, Mishra, 2021 --- simplified asymptotic form]\label{th:drlm}
Let $d, s, k \in \N$ with $k < s$ and let $f \in W^{s,\infty}\big((0,1)^d\big)$. For every $N$
large enough there is a network $\Phi$ with activation $\tanh$, \emph{two} hidden layers and width
$O(N^d)$, whose weights are bounded by an explicit power of $N$, such that
\[
\norm{f - \Phi}_{W^{k,\infty}((0,1)^d)} \;\le\; C(d,s,k,f)\ \frac{(\log N)^{k}}{N^{\,s-k}} .
\]
\end{theorem}

\begin{theorem}[Gühring, Kutyniok, Petersen, 2020 --- simplified asymptotic form]\label{th:gkp}
Let $1 \le p \le \infty$, let $n \ge 2$ be the Sobolev regularity of the target and let
$s \in [0,1]$ be the order in which the error is measured. For every $\varepsilon \in (0,\tfrac12)$
and every $f$ in the unit ball of $W^{n,p}((0,1)^d)$ there is a ReLU network $\Phi$ of depth
$O\big(\log(1/\varepsilon)\big)$ with
\[
O\Big( \varepsilon^{-d/(n-s)}\,\log^2(1/\varepsilon) \Big) \ \text{ weights, such that }\quad
\norm{f - \Phi}_{W^{s,p}((0,1)^d)} \le \varepsilon .
\]
\end{theorem}

\begin{remark}[Read the sources, not these boxes]\label{rem:sobolev-caveat}
Both statements above are \emph{simplified asymptotic forms}, given for fixed data, and neither is
the literal theorem of the paper it is attributed to. In De Ryck--Lanthaler--Mishra (Theorem 5.1) the
constant depends on $f$ in a way that does not factor out as $\norm{f}_{W^{s,\infty}}$, the
logarithm appears inside a factor of the form $\max\{R^k,\ \log^k(\beta N^{s+d+2})\}$, and the
conclusion holds only beyond an explicit threshold $N_0(d)$. In Gühring--Kutyniok--Petersen the
relevant statements are Theorem 4.1 and Corollary 4.2, whose hypotheses on the network class and on
$(n,s,p)$ are more restrictive than the shorthand above suggests --- note in particular that the
error order $s$ is a real number in $[0,1]$, not an integer, and that the regularity is $n \ge 2$.
Anyone using either result quantitatively must go back to the paper; these boxes are here to convey
the shape of the answer, not its constants.
\end{remark}

\begin{keybox}
\textbf{Three readings.}
\begin{enumerate}[label=(\alph*)]
\item Theorem \ref{th:gkp} is \emph{restricted to error orders $s \le 1$}: ReLU realisations do not
in general belong to $W^{2,p}$, their second derivatives being distributions or measures rather than
$L^p$ functions. This is the same obstruction as in Proposition \ref{prop:relu}, seen from the side
of rates.
\item The rates are of the form $N^{-(s-k)}$ with width $O(N^d)$: in other words,
\emph{the curse of dimensionality is very much present} for Sobolev classes. Barron's theorem
circumvents it only by changing the class of functions.
\item Differentiation costs. Within this family of bounds, raising by one the order in which the
error is measured loses a whole order in the rate; for a second-order PDE one therefore loses two
orders relative to what the same architecture achieves in $L^\infty$. Keep the logical status
straight: these are existential \emph{upper} bounds, so what the statement records is a limitation
of what is currently provable, not a proven lower bound, and certainly not a complete causal account
of observed numerical performance. It is nevertheless a feature shared by the constructions we know,
and it is a reasonable first thing to have in mind when comparing PINNs with classical methods on
smooth problems in low dimension.
\end{enumerate}
\end{keybox}

% ================================================================
\section{Depth versus width}

The theorems above use one or two hidden layers. Does depth bring anything? The answer, due to
Telgarsky, is yes, and exponentially so.

\begin{definition}[Hat map and sawtooth]
Let $\Delta : [0,1] \to [0,1]$ be defined by $\Delta(x) = 2x$ for $x \le \tfrac12$ and
$\Delta(x) = 2(1-x)$ for $x \ge \tfrac12$. Write
$\Delta^{\circ k} = \Delta \circ \cdots \circ \Delta$ ($k$ times).
\end{definition}

\begin{center}
\begin{tikzpicture}[scale=1]
\begin{axis}[
  width=13cm, height=4cm,
  xmin=0, xmax=1, ymin=0, ymax=1.08,
  axis lines=left, tick align=outside,
  xtick={0,0.25,0.5,0.75,1}, ytick={0,1},
  xlabel style={font=\scriptsize}, ylabel style={font=\scriptsize},
  tick label style={font=\scriptsize, color=coursgrey},
  axis line style={color=coursgrey},
  legend style={font=\scriptsize, draw=none, fill=none, at={(1.0,1.25)}, anchor=north east},
  legend cell align=left
]
\addplot[coursteal, thick, mark=none] coordinates {(0,0)(0.5,1)(1,0)};
\addlegendentry{$\Delta$}
\addplot[courswine, thick, mark=none] coordinates {(0,0)(0.25,1)(0.5,0)(0.75,1)(1,0)};
\addlegendentry{$\Delta^{\circ 2}$}
\addplot[coursochre, thick, mark=none, densely dashed] coordinates
  {(0,0)(0.125,1)(0.25,0)(0.375,1)(0.5,0)(0.625,1)(0.75,0)(0.875,1)(1,0)};
\addlegendentry{$\Delta^{\circ 3}$}
\end{axis}
\end{tikzpicture}
\end{center}

\begin{lemma}[Deep realisation]\label{lem:saw-deep}
One has $\Delta(x) = 2\ReLU(x) - 4\ReLU(x - \tfrac12)$ for $x \in [0,1]$. Consequently
$\Delta^{\circ k}$ is realised by a ReLU network of depth $k+1$ and width $2$, hence with $2k$
hidden neurons in total.
\end{lemma}

\begin{proof}
For $x \in [0,\tfrac12]$: $2x - 0 = 2x = \Delta(x)$. For $x \in [\tfrac12,1]$:
$2x - 4(x - \tfrac12) = 2 - 2x = \Delta(x)$. Composing $k$ blocks, each consisting of a hidden layer
of two ReLU neurons followed by a linear combination, gives a network of depth $k+1$ (that is, $k$
hidden layers) and width $2$.
\end{proof}

\begin{lemma}[Counting pieces]\label{lem:saw-pieces}
\begin{enumerate}[label=(\roman*)]
\item $\Delta^{\circ k}$ is piecewise affine on $[0,1]$ with exactly $2^k$ pieces: on each interval
$I_{k,j} = [\,j\,2^{-k}, (j+1)2^{-k}\,]$, $0 \le j < 2^k$, it is affine and sweeps out all of
$[0,1]$.
\item Every function $\R \to \R$ of the form
$x \mapsto c_0 + \sum_{i=1}^n c_i \ReLU(w_i x + b_i)$ is piecewise affine with at most $n+1$ pieces.
\end{enumerate}
\end{lemma}

\begin{proof}
(i) Induction on $k$. For $k=1$ this is the definition. Assume the property at rank $k$: on
$I_{k,j}$, $\Delta^{\circ k}$ is affine and bijective from $I_{k,j}$ onto $[0,1]$. Then
$\Delta^{\circ(k+1)} = \Delta \circ \Delta^{\circ k}$ restricted to $I_{k,j}$ is the composition of
an affine bijection $I_{k,j} \to [0,1]$ with $\Delta$, hence is affine on each half of $I_{k,j}$ and
sweeps out $[0,1]$ on each. This gives $2\cdot 2^k = 2^{k+1}$ pieces, which are exactly the
$I_{k+1,j'}$.

(ii) Each term $\ReLU(w_i x + b_i)$ is affine on $(-\infty, -b_i/w_i]$ and on
$[-b_i/w_i, +\infty)$ (and affine everywhere if $w_i = 0$). The sum is therefore affine on each
connected component of the complement of the $n$ points $\{-b_i/w_i\}$, of which there are at most
$n+1$.
\end{proof}

\begin{theorem}[Depth/width separation, after Telgarsky 2016]\label{th:telgarsky}
For every $k \ge 1$ the function $\Delta^{\circ k}$ is exactly realised by a ReLU network of depth
$k+1$ with $2k$ neurons, whereas any ReLU network with a \emph{single} hidden layer realising it
exactly has at least $2^k - 1$ neurons.
\end{theorem}

\begin{proof}
The first assertion is Lemma \ref{lem:saw-deep}. For the second: if a network with one hidden layer
of $n$ neurons realises $\Delta^{\circ k}$ on $[0,1]$, then by Lemma \ref{lem:saw-pieces}\,(ii) that
function has at most $n+1$ affine pieces, whereas it has exactly $2^k$ by (i). Hence
$n + 1 \ge 2^k$, that is, $n \ge 2^k - 1$.
\end{proof}

\begin{corollary}[Piece counting at arbitrary depth]\label{cor:pieces}
A ReLU network $\R \to \R$ of depth $L$ whose hidden layers have $n_1,\dots,n_{L-1}$ neurons is
piecewise affine with at most $\prod_{\ell=1}^{L-1}(n_\ell + 1)$ pieces. Hence, at a fixed total
budget $n = \sum_\ell n_\ell$ of neurons, the number of pieces available grows exponentially in the
depth and only linearly in the size of a single hidden layer.
\end{corollary}

\begin{proof}
Exercise \ref{ex:telgarsky}, item 4.
\end{proof}

\begin{remark}[What Telgarsky actually proves]\label{rem:telgarsky-real}
Theorem \ref{th:telgarsky} is a statement about \emph{exact representation}, obtained by counting
pieces; it is the elementary version of the phenomenon and it is the one we prove. Telgarsky's
published result is stronger and different in shape: he exhibits a function computed by a ReLU
network with $\Theta(k^3)$ layers and $\Theta(1)$ nodes per layer such that every network with
$O(k)$ layers built from semi-algebraic gates (ReLU, maximum, indicator, piecewise polynomial)
needs $\Omega(2^k)$ nodes to come within a fixed constant of it in $L^1$. So the separation survives
when one merely asks for \emph{approximation}, not exact representation. We do not prove this here;
see the COLT 2016 paper.
\end{remark}

\begin{remark}[Relevance for PINNs]
The lesson is that depth is useful for representing \emph{highly oscillatory} functions with few
parameters. In practice PINNs use moderately deep (4 to 8 layers) and narrow (20 to 100 neurons)
networks: a compromise between expressive power and the conditioning of the optimisation, which we
shall analyse in Session 6.
\end{remark}

% ================================================================
\section{Summary: what to retain}

\begin{keybox}
\begin{enumerate}[label=\textbf{\arabic*.}, leftmargin=*]
\item \textbf{The class is not linear.} Fixed-size realisation classes are in general nonlinear and
nonconvex; Proposition \ref{prop:nonconvex} exhibits this already for $\Sigma_1(\sigma)$, and
Remark \ref{rem:nonclosed} adds that they need not be closed. Céa's lemma and the Galerkin machinery
no longer apply, and uniqueness and attainment of the minimiser are no longer automatic. This is one
major source of the difficulties in this course --- PDE stability, quadrature, sampling and the
optimiser each contribute difficulties of their own, and we shall meet them in turn.
\item \textbf{Universality holds, and by itself is useless.} For every continuous nonpolynomial
$\sigma$, $\Sigma(\sigma)$ is dense in $C(K)$ (Theorems \ref{th:cybenko} and \ref{th:llps}). The
proof is nonconstructive and gives no rate.
\item \textbf{Rates are decided by the target class at least as much as by the architecture.}
Activation, depth, width and weight constraints all enter the bounds; but the qualitative behaviour
--- whether the curse of dimensionality appears at all --- is set by the class the target is assumed
to lie in. On a Sobolev class the curse is present (Theorems \ref{th:drlm}, \ref{th:gkp}); on the
Barron class it disappears (Theorem \ref{th:barron}), at the price of a spectral hypothesis whose
constant may itself grow with $d$.
\item \textbf{Approximating derivatives costs orders.} For an operator of order $m$ it is the
$W^{m,\infty}$ norm that matters, and each order of differentiation costs one order of rate. The
regularity of $\sigma$ is constrained by the order of the PDE (Proposition \ref{prop:relu}).
\item \textbf{All of this is existential.} None of these results says that an algorithm will find
the right $\theta$. That question is formalised in Session 4 (error decomposition) and treated in
Session 6 (analysis of training) --- where it will turn out that this is where a substantial part of
the practical difficulty lies.
\end{enumerate}
\end{keybox}

\medskip
\noindent
\textbf{Session 2.} How does one compute $\nabla_\theta J$ and $\partial^\alpha_x \Phi_\theta$?
Automatic differentiation in forward and reverse mode; nonconvex optimisation; what Adam actually
does, and why the PINN literature systematically chains Adam and L-BFGS.

% ================================================================
\newpage
\section{Problem set}

\begin{exercise}[Regularity of the activation and the strong residual]\label{ex:relu}
\leavevmode
\begin{enumerate}[label=\textbf{\arabic*.}]
\item Let $\Phi(x) = \sum_{i=1}^n c_i \ReLU(w_i x + b_i) + c_0$ with $d = 1$. Show that $\Phi$ is
continuous and piecewise affine, and compute $\Phi''$ in the sense of distributions.
\item Deduce that for the problem $-u'' = f$ on $(0,1)$ the strong residual $\Rr(v) = -v'' - f$ is
not a function when $\sigma = \ReLU$, so that
$\theta \mapsto \frac1N \sum_{i=1}^N \abs{\Rr(\Phi_\theta)(x_i)}^2$ is not defined by the PDE. Then
ask the more useful question: what does an automatic-differentiation library actually compute here,
and what is the resulting loss equal to?
\item Now let $\sigma \in C^m(\R)$ and let $\Dd$ be a differential operator of order $m$ with
continuous coefficients. Show that $\Rr(\Phi_\theta)$ is continuous, hence that pointwise evaluation
is legitimate.
\item \emph{(Weak formulation.)} Consider the Deep Ritz functional
$\Ee(v) = \int_0^1 \big( \tfrac12 \abs{v'}^2 - f v \big)$. For which regularity of $\sigma$ is it
defined? Comment.
\end{enumerate}
\end{exercise}

\begin{exercise}[Structure of the class]\label{ex:structure}
Let $\sigma = \tanh$ and $d \ge 2$.
\begin{enumerate}[label=\textbf{\arabic*.}]
\item Show that $\Sigma_n(\sigma)$ is a cone: stable under multiplication by a scalar.
\item Show that $\Sigma_n(\sigma) + \Sigma_m(\sigma) \subseteq \Sigma_{n+m}(\sigma)$, and that the
inclusion is strict for $n = m = 1$.
\item Revisit the proof of Proposition \ref{prop:nonconvex} and adapt it to the case $d = 1$: show
that $\tfrac12\big(\tanh(x) + \tanh(2x)\big) \notin \Sigma_1(\tanh)$.
\emph{Hint: counting inflection points is the natural first attempt --- show that it fails ---
then use oddness, the limits at $\pm\infty$, and the Taylor expansions at $0$.}
\item Deduce that the functional $\theta \mapsto \norm{f - \Phi_\theta}_{L^2}^2$ can fail to be
convex in $\theta$ even when $f$ is analytic and exactly representable in the class.
\end{enumerate}
\end{exercise}

\begin{exercise}[The polynomial case]\label{ex:poly}
\leavevmode
\begin{enumerate}[label=\textbf{\arabic*.}]
\item Let $\sigma(t) = t^2$. Describe $\Sigma(\sigma)$ exactly for $d = 1$, then for general $d$.
\item Show that with $\sigma(t) = t^2$ and \emph{fixed} depth $L$, one has
$\Nn^\sigma_{L,W}(d) \subseteq \Pp_{2^{L-1}}(\R^d)$ for every width $W$, and deduce that this class
is not dense in $C(K)$.
\item Now let the depth vary. Using $xy = \tfrac14\big((x+y)^2 - (x-y)^2\big)$, show that
$\bigcup_{L,W} \Nn^{t^2}_{L,W}(d)$ contains every polynomial, and is therefore dense in $C(K)$.
Reconcile this with item 2.
\item In what way does this example illuminate the statement of Theorem \ref{th:llps}? Which class
does that theorem quantify over?
\item Let $\sigma$ be $C^\infty$, nonpolynomial, but with $\sigma^{(k)}(0) = 0$ for every $k \ge 3$.
Does such a $\sigma$ exist? What happens to the proof of
Proposition \ref{prop:llps-smooth}?
\end{enumerate}
\end{exercise}

\begin{exercise}[Maurey's lemma and its optimality]\label{ex:maurey}
\leavevmode
\begin{enumerate}[label=\textbf{\arabic*.}]
\item Reprove Lemma \ref{lem:maurey}, writing out the cancellation of the cross terms in detail.
\item Let $\Hh = \ell^2$ and $G = \{\pm e_j : j \ge 1\}$ (a Hilbert basis). Let
$f = \frac1m \sum_{j=1}^m e_j$. Check that $f \in \co(G)$ and compute $\norm f$. What does the lemma
give?
\item Show that in this example, with $m = 2n$, any average of $n$ elements of $G$ is at distance
$\ge \tfrac12 n^{-1/2}$ from $f$. Deduce that the rate $n^{-1/2}$ is optimal.
\item \emph{(Open-ended.)} Is the rate $n^{-1/2}$ improved if $\Hh$ has finite dimension $D$?
Compare with $D^{-1/2}$ and comment on the role of $D$ in applications.
\end{enumerate}
\end{exercise}

\begin{exercise}[The Barron class]\label{ex:barron}
\leavevmode
\begin{enumerate}[label=\textbf{\arabic*.}]
\item Prove item (ii) of Proposition \ref{prop:barron-elem}: $f \in \Gamma_C$ implies $f$ Lipschitz
with constant $\le C$.
\item Give an upper bound for $C_f$ when $f(x) = \ee^{-\abs x^2/2}$ in dimension $d$, and an
equivalent of that bound as $d \to \infty$. Why is one not entitled to write an equality here?
\item Compute $C_f$ \emph{exactly} when $f(x) = \sin(K x_1)$, $K > 0$: an upper bound by exhibiting
a representing measure, and a matching lower bound from item (ii) of
Proposition \ref{prop:barron-elem}. Comment: is the Barron class a regularity class? Is it
indifferent to regularity?
\item Let $f(x) = \prod_{i=1}^d \cos(x_i)$. Show that $C_f \le \sqrt d$ and discuss: is Barron's
theorem useful for this function in high dimension?
\item \emph{(Link with Session 6.)} Is a PDE solution belonging to $\Gamma_C$ with $C$ large hard
for a PINN? Formulate a conjecture and keep it in mind for Session 6.
\end{enumerate}
\end{exercise}

\begin{exercise}[Telgarsky's sawtooth]\label{ex:telgarsky}
\leavevmode
\begin{enumerate}[label=\textbf{\arabic*.}]
\item Check that $\Delta(x) = 2\ReLU(x) - 4\ReLU(x - \tfrac12)$ on $[0,1]$.
\item Prove by induction that $\Delta^{\circ k}$ has exactly $2^k$ affine pieces.
\item Deduce Theorem \ref{th:telgarsky}.
\item \emph{(Generalisation.)} Show that a ReLU network $\R \to \R$ of depth $L$ whose hidden layers
have $n_1,\dots,n_{L-1}$ neurons has at most $\prod_{\ell}(n_\ell + 1)$ affine pieces. Compare with
the number of neurons $\sum_\ell n_\ell$ and conclude as to the exponential advantage of depth.
\item What becomes of this advantage if ReLU is replaced by $\tanh$?
\end{enumerate}
\end{exercise}

% ================================================================

\ifsolutions
\newpage
\section{Solutions}

\subsection*{Exercise \ref{ex:relu}}

\textbf{1.} Each $\ReLU(w_i x + b_i)$ is affine on $(-\infty, \tau_i]$ and on $[\tau_i,+\infty)$
with $\tau_i = -b_i/w_i$ (if $w_i \ne 0$; otherwise the function is constant). Hence $\Phi$ is
affine on each component of the complement of $\{\tau_1,\dots,\tau_n\}$, and continuous by
continuity of $\ReLU$.

In the sense of distributions, for $w \ne 0$,
\[
\big(\ReLU(w x + b)\big)'' = \abs{w}\,\delta_{-b/w}.
\]
The quickest check is the jump of the first derivative: $\big(\ReLU(wx+b)\big)' = w\,H(wx+b)$ with
$H$ the Heaviside function, which for $w > 0$ jumps from $0$ to $w$ at $\tau = -b/w$, and for
$w < 0$ jumps from $w$ to $0$; in both cases the jump equals $\abs{w}$. Equivalently, in the change
of variable $y = wx+b$ the two differentiations contribute $w^2$ while the Jacobian contributes
$\abs{w}^{-1}$, and $w^2/\abs{w} = \abs{w}$. Therefore
\[
\Phi'' = \sum_{i \,:\, w_i \ne 0} c_i\, \abs{w_i}\, \delta_{\tau_i},
\]
a finite combination of Dirac masses. \emph{(Note the modulus. Writing $w_i^2$ here --- keeping the
two differentiations and forgetting the Jacobian --- is the standard slip.)}

\textbf{2.} $\Rr(\Phi_\theta) = -\Phi_\theta'' - f$ is the sum of a singular measure and a function,
so it is not a function and $\Rr(\Phi_\theta)(x_i)$ is not defined by the PDE. But the sharp
statement --- the one to give in class --- is about what happens in code. Automatic differentiation
does not return the distributional derivative; it returns the classical one, which exists and equals
$0$ at every point off the finite set $\{\tau_i\}$. At generic collocation points the interior loss
is therefore numerically well defined and equal to $\frac1N\sum_i \abs{f(x_i)}^2$: a constant, with
identically zero gradient in $\theta$. So the failure mode is not ``the code refuses to run'' but
``the residual term is blind'' --- it contains no information about $\theta$ whatsoever. Saying only
that the loss is meaningless hides this distinction, which is the one a student will meet in
practice.

\textbf{3.} By Proposition \ref{prop:stable}\,(iv), $\Phi_\theta \in C^m$; the derivatives
$\partial^\alpha \Phi_\theta$ for $\abs\alpha \le m$ are continuous, and $\Dd[\Phi_\theta]$, a
finite sum of products of continuous functions, is continuous too. Pointwise evaluation is therefore
well defined.

\textbf{4.} $\Ee(v)$ requires $v \in H^1(0,1)$. A \emph{locally Lipschitz} activation suffices on a
bounded domain --- $\ReLU$ qualifies --- but mere continuity of $\sigma$ does not: a continuous
realisation need not lie in $H^1$. This is precisely the point of variational formulations: they
restore compatibility with nonsmooth activations and lower the order of differentiation that
automatic differentiation must handle (Session 7). In exchange they introduce a quadrature of an
energy integral, whose variance becomes the new problem. Note also that VPINNs, which integrate by
parts against test functions, impose a different requirement again; Deep Ritz and VPINN should not
be described by a single condition on $\sigma$.

\subsection*{Exercise \ref{ex:structure}}

\textbf{1.} $\lambda \sum_i c_i \sigma(w_i\cdot x + b_i) = \sum_i (\lambda c_i)\sigma(w_i \cdot x+b_i)$.

\textbf{2.} The inclusion is concatenation of the sums. For strictness with $n = m = 1$:
Proposition \ref{prop:nonconvex} provides $f_1 + f_2 \in \Sigma_2 \setminus \Sigma_1$ (up to the
constant $\tfrac12$, absorbed by item 1).

\textbf{3.} We must show that $g(x) = \tfrac12(\tanh x + \tanh 2x)$ is not of the form
$c\tanh(wx+b)$. Every function $c\tanh(wx+b)$ with $cw \ne 0$ is strictly monotone and has
\emph{exactly one} inflection point, at $x = -b/w$ (because $\tanh''(t) = -2\tanh t\,\sech^2 t$
vanishes only at $t = 0$, where it changes sign). But
\[
g''(x) = \tfrac12\Big( -2\tanh(x)\sech^2(x) - 8 \tanh(2x)\sech^2(2x) \Big) = -\big(\psi(x) + 4\psi(2x)\big),
\qquad \psi(t) := \tanh(t)\sech^2(t).
\]
The function $\psi$ is odd, strictly positive on $(0,+\infty)$ and strictly negative on
$(-\infty,0)$. Hence $\psi(x) + 4\psi(2x) > 0$ for $x > 0$ and $< 0$ for $x < 0$: $g''$ vanishes
only at $0$. So counting inflection points yields \emph{no} contradiction --- a further argument is
needed.

Compare the asymptotic behaviours instead. If $g(x) = c\tanh(wx+b)$, then
$\lim_{x\to+\infty} g = c\,\mathrm{sgn}(w)$ and $\lim_{x \to -\infty} g = -c\,\mathrm{sgn}(w)$,
whereas $g(x) \to 1$ and $g(x) \to -1$. Hence $c\,\mathrm{sgn}(w) = 1$, and after adjusting signs we
may assume $c = 1$, $w > 0$. Moreover $g$ is odd (as $\tanh$ is), so
$\tanh(wx+b) = -\tanh(-wx+b) = \tanh(wx - b)$ for every $x$, which forces $b = 0$ by injectivity of
$\tanh$. It remains to rule out $g(x) = \tanh(wx)$ for some $w > 0$. Now
\[
g'(0) = \tfrac12(1 + 2) = \tfrac32, \qquad \big(\tanh(w\cdot)\big)'(0) = w \ \Longrightarrow\ w = \tfrac32,
\]
but then, expanding to order $3$ at $0$ (using $\tanh t = t - t^3/3 + O(t^5)$):
\[
g(x) = \tfrac32 x - \tfrac12\big(\tfrac{1}{3} + \tfrac{8}{3}\big) x^3 + O(x^5) = \tfrac32 x - \tfrac32 x^3 + O(x^5),
\qquad
\tanh(\tfrac32 x) = \tfrac32 x - \tfrac{9}{8}x^3 + O(x^5),
\]
and $\tfrac32 \ne \tfrac98$. Hence $g \notin \Sigma_1(\tanh)$.

\textbf{4.} There is a two-line counterexample, and it needs nothing from item 3. Take $d = 1$,
$\Omega = (-1,1)$ (any set of positive measure will do), $f = \tanh$, and the two parameters
\[
\theta_+ = (c,w,b) = (1,1,0), \qquad \theta_- = (c,w,b) = (-1,-1,0).
\]
Since $\tanh$ is odd, $\Phi_{\theta_+}(x) = \tanh(x)$ and
$\Phi_{\theta_-}(x) = -\tanh(-x) = \tanh(x)$: the two parameters realise the \emph{same} function,
so
\[
J(\theta_+) = J(\theta_-) = 0 .
\]
But $\tfrac12(\theta_+ + \theta_-) = (0,0,0)$, whose realisation is the zero function, whence
\[
J\Big( \tfrac{\theta_+ + \theta_-}{2} \Big) = \norm{\tanh}_{L^2(\Omega)}^2 > 0 .
\]
A convex function cannot exceed the larger of its values at the endpoints of a segment; hence $J$ is
not convex in $\theta$ --- and this for a target $f$ that is analytic and lies exactly in
$\Sigma_1(\tanh)$, so that the global minimum is $0$ and is attained.

\emph{A warning about the argument one is tempted to give instead.} It is not enough to observe that
$\theta \mapsto \Phi_\theta$ is not affine. Composing a convex functional with a nonaffine map can
perfectly well yield a convex function: $t \mapsto t^2$ composed with $x \mapsto x^2$ gives
$x \mapsto x^4$, which is convex. Nonlinearity of the parametrisation is the \emph{reason to
expect} nonconvexity; it is not a proof of it. The example above is a proof.

\subsection*{Exercise \ref{ex:poly}}

\textbf{1.} For $d = 1$: $\sigma(wx+b) = (wx+b)^2 = w^2x^2 + 2wbx + b^2$, so
$\Sigma(\sigma) = \Pp_2(\R)$ --- the polynomials of degree $\le 2$ (one does obtain all three
monomials: $b=0$ gives $x^2$, $w=0$ gives the constants, and a combination gives $x$). For general
$d$: $\Sigma(\sigma) = \Pp_2(\R^d)$, by the same argument together with the polarisation of Step 2
in Proposition \ref{prop:llps-smooth} applied with $k = 2$.

\textbf{2.} By induction: one hidden layer maps a polynomial of degree $\le D$ to a polynomial of
degree $\le 2D$ (composition with $t \mapsto t^2$), and affine combinations preserve the degree.
Starting from $D = 1$ (the affine map $A_1$), after $L-1$ hidden layers one has $D \le 2^{L-1}$.
Hence $\Nn^{t^2}_{L,W}(d) \subseteq \Pp_{2^{L-1}}(\R^d)$ for every width $W$; the right-hand side is
finite-dimensional, hence closed and proper in $C(K)$, so this class is not dense --- \emph{at that
fixed depth}.

\textbf{3.} Let the depth grow. The identity $xy = \tfrac14\big((x+y)^2 - (x-y)^2\big)$ shows that
a single hidden layer of two quadratic units followed by a linear combination computes the product
of any two quantities already available at that layer. Starting from the inputs
$x_1,\dots,x_d$ (and the constant $1$, available as $\sigma(0\cdot x + 1)$), one builds every
degree-2 monomial with one hidden layer, every monomial of degree $\le 2^m$ with $m$ layers of
pairwise products, and any linear combination of them in the output layer. Hence
\[
\bigcup_{L,W} \Nn^{t^2}_{L,W}(d) \ \supseteq\ \Pp(\R^d),
\]
which is dense in $C(K)$ by Stone--Weierstrass.

There is no contradiction with item 2: density fails at each \emph{fixed} depth and holds only for
the union over all depths. The correct slogan is therefore \emph{a polynomial activation is not
universal at fixed depth}. The tempting slogan --- ``a quadratic activation never gets you out of
the polynomials, so nothing can help'' --- is false as a conclusion: it never does get you out of
the polynomials, and that is precisely enough, because the polynomials are dense.

\textbf{4.} Theorem \ref{th:llps} quantifies over $\Sigma(\sigma)$, the span of a \emph{single}
hidden layer. For that class, nonpolynomiality of $\sigma$ is necessary and sufficient, and
$\sigma = t^2$ is not universal. Item 3 shows that the analogous statement for the union over all
depths is \emph{false}. Keeping track of which class a theorem quantifies over is not pedantry
here: it changes the answer. (The proof of Proposition \ref{prop:llps-smooth} makes the mechanism
visible: it never composes layers, it only differentiates in the inner scaling parameter of one
layer.)

\textbf{5.} Yes, such a $\sigma$ exists: take $\sigma(t) = \ee^{-1/t^2}$ for $t > 0$ and
$\sigma(t) = 0$ for $t \le 0$; all its derivatives vanish at $0$, and it is not polynomial. The
proof of Proposition \ref{prop:llps-smooth} does not use $b = 0$ but the existence, for each $k$, of
some $b_k$ with $\sigma^{(k)}(b_k) \ne 0$: here it suffices to take $b_k > 0$. The argument
therefore goes through unchanged. This is a good comprehension check: nonpolynomiality is a
\emph{global} property, not a local one at a point --- a function all of whose derivatives vanish at
one point can still be universal.

\subsection*{Exercise \ref{ex:maurey}}

\textbf{1.} See the proof of Lemma \ref{lem:maurey}. In detail, for the cross terms:
\[
\E \Big\| \sum_{i=1}^n (G_i - f)\Big\|^2
= \sum_{i,j} \E \scal{G_i - f}{G_j - f}
= \sum_i \E \norm{G_i - f}^2 + \sum_{i \ne j} \scal{\E[G_i] - f}{\E[G_j] - f},
\]
and the second term vanishes because $\E[G_i] = f$ and $G_i, G_j$ are independent for $i \ne j$.

\textbf{2.} $f = \frac1m\sum_{j\le m} e_j$ is indeed a convex combination of the $e_j \in G$, and
$\norm f^2 = m \cdot (1/m)^2 = 1/m$. Here $\beta = 1$, so the lemma gives the existence of
$g_1,\dots,g_n \in G$ with $\norm{f - \frac1n\sum g_i}^2 \le (1 - 1/m)/n$.

\textbf{3.} Let $S = \frac1n \sum_{i=1}^n g_i$ with $g_i \in \{\pm e_j\}$. Then $S$ is supported on
at most $n$ coordinates. Let $J$ be the set (of cardinality $\le n$) of those coordinates. Then
\[
\norm{f - S}^2 \ge \sum_{j \le m,\ j \notin J} \frac{1}{m^2} \ge \frac{m - n}{m^2} = \frac1m - \frac{n}{m^2}.
\]
Choosing $m = 2n$ gives $\norm{f-S}^2 \ge \frac{1}{2n} - \frac{1}{4n} = \frac{1}{4n}$, that is,
$\norm{f - S} \ge \tfrac12 n^{-1/2}$. The rate $n^{-1/2}$ of the lemma is therefore optimal up to a
constant.

\textbf{4.} The tempting answer --- ``in dimension $D$ one can be exact as soon as $n \ge D$'' ---
is false, and instructively so. Lemma \ref{lem:maurey} produces an \emph{equally weighted} average
$\frac1n\sum_i g_i$, and a convex combination with, say, irrational weights is in general not an
equal-weight average of elements of $G$, in any dimension whatsoever.

What finite dimension does buy is the following. By Carathéodory's theorem, any $f \in \co(G)$ with
$G \subseteq \Hh$ and $\dim \Hh = D$ is a convex combination $\sum_{j=1}^{D+1}\lambda_j h_j$ of at
most $D+1$ elements of $G$. Round the weights to multiples of $1/n$: choose integers $k_j \ge 0$ with
$\sum_j k_j = n$ and $\abs{\lambda_j - k_j/n} \le 1/n$ for each $j$. Then
\[
\Big\| f - \frac1n\sum_{j} k_j h_j \Big\|
\ \le\ \sum_{j=1}^{D+1} \Big|\lambda_j - \frac{k_j}{n}\Big|\,\norm{h_j}
\ \le\ \frac{(D+1)\,\beta}{n},
\]
an $O(D\beta/n)$ rate, which beats $\beta/\sqrt n$ once $n \gtrsim D^2$. So the dimension enters
through a crossover point, not by making the problem exactly solvable.

In applications $\Hh = L^2(\mu)$ is infinite-dimensional, but the \emph{effective} dimension of the
problem --- the decay of the eigenvalues of the covariance operator --- plays an analogous role.
This is the point of view we shall meet again in Session 9, on Kolmogorov $n$-widths in operator
learning.

\subsection*{Exercise \ref{ex:barron}}

\textbf{1.} See Proposition \ref{prop:barron-elem}\,(ii). The key point is
$\abs{\ee^{\ii a} - \ee^{\ii b}} = 2\abs{\sin\frac{a-b}{2}} \le \abs{a - b}$.

\textbf{2.} Take $\dd\nu(\xi) = (2\pi)^{-d/2}\ee^{-\abs\xi^2/2}\dd\xi$ (check by the inversion
formula that it represents $f$ on all of $\R^d$). Hence
\[
C_f \ \le\ \int_{\R^d}\abs\xi\,(2\pi)^{-d/2}\ee^{-\abs\xi^2/2}\dd\xi = \E\abs{Z},\quad Z\sim\Nn(0,I_d),
\qquad
\E\abs Z = \sqrt2\,\frac{\Gamma\big(\frac{d+1}2\big)}{\Gamma\big(\frac d2\big)} .
\]
By Stirling's formula, $\E\abs Z = \sqrt d\,\big(1 - \tfrac{1}{4d} + O(d^{-2})\big) \sim \sqrt d$.
The growth is therefore \emph{polynomial and very slow}: for this structured family, Barron's bound
avoids the generic Sobolev curse in high dimension. (It is a comparison of a bound for a specific
function with a bound uniform over a Sobolev ball, so it is not a like-for-like contest.)

This is the answer to the second half of the question. We have exhibited \emph{one} representing
measure; a cheaper one, agreeing with $f$ only on $B_r$ and differing from it outside, is not
excluded, and the infimum need not be attained
(Remark \ref{rem:barron-hygiene}\,(c)). Writing $C_f = \E\abs Z$ would therefore be an unproved
claim. An upper bound is in any case all that Theorem \ref{th:barron} consumes, since its conclusion
is increasing in $C_f$.

\textbf{3.} Take $\nu = \frac{1}{2\ii}\big(\delta_{Ke_1} - \delta_{-Ke_1}\big)$, a finite complex
measure with $\int \ee^{\ii \xi\cdot x}\dd\nu(\xi) = \sin(Kx_1)$ for every $x$. Its total variation
is $\abs\nu = \tfrac12(\delta_{Ke_1} + \delta_{-Ke_1})$, so
$\int\abs\xi\dd\abs\nu = \tfrac12 K + \tfrac12 K = K$, whence $C_f \le K$.

Here, unlike in item 2, the matching lower bound is immediate and should be given: by
Proposition \ref{prop:barron-elem}\,(ii), $C_f$ dominates the Lipschitz constant of $f$ on $B_r$,
and $\nabla f(x) = K\cos(Kx_1)e_1$ gives $\abs{\nabla f(0)} = K$. Hence $C_f \ge K$ and
\[
C_{\sin(K x_1)} = K .
\]
Without this lower bound the phrase ``the seminorm grows linearly with frequency'' would be an
unproved reading of an upper bound.

This is also the example that forces the definition to use measures rather than $L^1$ densities:
$\sin(Kx_1)$ has no $L^1$ Fourier density, and mollifying the two atoms would destroy the
\emph{exact} identity on $B_r$ that the definition requires.

The moral, stated carefully. The spectral Barron seminorm is not determined by classical smoothness
order alone: even an analytic function can have a seminorm that grows with frequency, so the class
is not \emph{graded} by smoothness. Nevertheless a finite seminorm does impose regularity, since the
representing measure has a finite first moment and the function therefore admits a $C^1$ extension
(Remark \ref{rem:barron-c1}). What one must not say is that a rough function with concentrated
spectrum is cheap in this class: within this definition there are no rough functions at all.

\textbf{4.} $f(x) = \prod_i \cos(x_i) = 2^{-d}\sum_{\epsilon \in \{\pm1\}^d} \ee^{\ii \epsilon\cdot x}$,
so take $\nu = 2^{-d}\sum_{\epsilon \in \{\pm1\}^d}\delta_\epsilon$: a positive measure carried by
the $2^d$ points $\epsilon \in \{\pm1\}^d$, each of mass $2^{-d}$, with $\abs\epsilon = \sqrt d$.
Hence $C_f \le 2^d \cdot 2^{-d}\cdot \sqrt d = \sqrt d$. Barron's theorem therefore says that
\[
n = O\!\left( \frac{r^2 d}{\varepsilon^2} \right) \quad \text{neurons \emph{suffice}}
\]
--- a sufficient width, not an equivalence and not a lower bound. The dependence on $d$ in this
sufficient width is \emph{linear}, to be compared with the $\varepsilon^{-d/s}$ of linear methods. This is a favourable
case and a nice example to give in class. (Note that $r$ itself typically grows like $\sqrt d$ if
the domain is $[-1,1]^d$, giving a final dependence in $d^2$ --- still polynomial.)

\textbf{5.} The expected conjecture: yes, with the logical status kept straight. High frequency
content makes the function expensive for \emph{approximation} --- a large Barron seminorm, hence a
potentially large \emph{sufficient} width according to Barron's bound, which is an upper bound and
not a necessity --- and, separately, hard for \emph{optimisation}. Session 6 will show that the second effect is by far the more severe --- but it
is important to state in which regime this is a theorem. In the regime where the analysis is
available (very wide networks near their initialisation, where the training dynamics are governed by
the neural tangent kernel and are effectively linear), the eigenvalues of that kernel associated
with high frequencies are far smaller than those associated with low ones, so gradient flow reduces
the high-frequency components of the error at a correspondingly slower rate. Outside that regime the
phenomenon is observed rather than proved. That gap between $\mathcal E_{\text{app}}$ and
$\mathcal E_{\text{opt}}$ is the real subject of this course.

\subsection*{Exercise \ref{ex:telgarsky}}

\textbf{1--3.} See Lemmas \ref{lem:saw-deep}, \ref{lem:saw-pieces} and Theorem
\ref{th:telgarsky}.

\textbf{4.} Write $p(\Phi)$ for the number of affine pieces of a continuous piecewise affine
function $\Phi : \R \to \R$. If $\Phi$ is piecewise affine with $p$ pieces and the next layer
applies $n$ ReLU units coordinatewise after an affine map, then each piece of $\Phi$ can be
subdivided into at most $n+1$ subpieces (each ReLU unit adds at most one breakpoint per piece).
Hence, layer after layer,
\[
p \le \prod_{\ell=1}^{L-1}(n_\ell + 1).
\]
For a fixed total number of neurons $n = \sum_\ell n_\ell$, this product is largest when the
$n_\ell$ are equal: $p \lesssim (n/(L-1) + 1)^{L-1}$, which grows \emph{exponentially} in $L$ and
only \emph{polynomially} in the width. A one-hidden-layer network achieves $p \le n+1$: this is the
exponential advantage of depth.

\textbf{5.} The counting argument collapses outright: $\tanh$ is analytic, so every $\tanh$ network
is analytic and has no ``pieces'' to count. Depth separation results do exist for smooth
activations, but they are obtained by other means --- Vapnik--Chervonenkis complexity, or
analyticity together with bounds on the number of zeros --- and we state no theorem about them here;
comparing their strength with Theorem \ref{th:telgarsky} would require fixing a common notion of
separation first, which we have not done. The honest conclusion is narrower and still worth drawing:
piece-counting is a tool specific to piecewise affine activations, which is one reason the theory of
ReLU networks is far more developed than that of $\tanh$ networks --- even though it is the latter
that PINNs use.

% ================================================================
\newpage
\fi

\section*{References for this session}
\addcontentsline{toc}{section}{References for this session}

\begin{itemize}[leftmargin=*, itemsep=4pt]
\item A.~Pinkus, \emph{Approximation theory of the MLP model in neural networks}, Acta Numerica
\textbf{8} (1999), 143--195. --- \emph{The reference for all the material in §§\,1--2.}
\item R.~DeVore, B.~Hanin, G.~Petrova, \emph{Neural network approximation}, Acta Numerica
\textbf{30} (2021), 327--444. --- \emph{A modern survey, very readable at MSc level.}
\item G.~Cybenko, \emph{Approximation by superpositions of a sigmoidal function}, Mathematics of
Control, Signals and Systems \textbf{2} (1989), 303--314.
\item M.~Leshno, V.~Y.~Lin, A.~Pinkus, S.~Schocken, \emph{Multilayer feedforward networks with a
nonpolynomial activation function can approximate any function}, Neural Networks \textbf{6} (1993),
861--867.
\item A.~R.~Barron, \emph{Universal approximation bounds for superpositions of a sigmoidal function},
IEEE Transactions on Information Theory \textbf{39} (1993), 930--945.
\item T.~De Ryck, S.~Lanthaler, S.~Mishra, \emph{On the approximation of functions by tanh neural
networks}, Neural Networks \textbf{143} (2021), 732--750.
\item I.~Gühring, G.~Kutyniok, P.~Petersen, \emph{Error bounds for approximations with deep ReLU
neural networks in $W^{s,p}$ norms}, Analysis and Applications \textbf{18} (2020), 803--859.
\item M.~Telgarsky, \emph{Benefits of depth in neural networks}, COLT 2016, 1517--1539.
\item P.~Petersen, M.~Raslan, F.~Voigtlaender, \emph{Topological properties of the set of functions
generated by neural networks of fixed size}, Foundations of Computational Mathematics \textbf{21}
(2021), 375--444. --- \emph{The reference for Remark \ref{rem:nonclosed}.}
\item S.~Mahan, E.~King, A.~Cloninger, \emph{Nonclosedness of sets of neural networks in Sobolev
spaces}, Neural Networks \textbf{137} (2021), 85--96.
\item P.~C.~Kainen, V.~K\r{u}rková, A.~Vogt, \emph{Approximation by neural networks is not
continuous}, Neurocomputing \textbf{29} (1999), 47--56. --- \emph{On the absence of a continuous
best-approximation selection; a different statement from non-closedness, with which it is often
conflated.}
\end{itemize}

\end{document}
